% --- SLIDE 5: DISTILLATION, SOFT TARGETS ---
\begin{frame}{Knowledge Distillation: A Small Student Learns from a Large Teacher}

    \begin{columns}[T]
        \begin{column}{0.5\textwidth}
            \begin{itemize}
                \small
                \item Train a large \textbf{teacher} for accuracy, then train a small \textbf{student} to
                      reproduce the teacher's \emph{output probabilities}, not just the labels (Hinton,
                      Vinyals \& Dean, 2015).
                \item A softmax with temperature $T$ softens those probabilities:
                      \[
                          q_i = \frac{\exp(z_i / T)}{\sum_j \exp(z_j / T)}
                      \]
                \item The soft targets carry what the teacher knows about \textbf{which wrong answers are
                      nearly right}. A hard label throws that away.
            \end{itemize}
        \end{column}
        \begin{column}{0.48\textwidth}
            \centering
            \setlength{\tabcolsep}{0.3em}
            \renewcommand{\arraystretch}{1.2}
            \begin{tabular}{l|cccc}
                & \footnotesize cow & \footnotesize car & \footnotesize dog & \footnotesize cat \\
                \hline
                \footnotesize Label          & $0$        & $1$    & $0$    & $0$ \\
                \footnotesize Teacher, $T{=}1$ & $10^{-6}$ & $0.9$  & $0.1$  & $10^{-9}$ \\
                \footnotesize Teacher, $T{=}5$ & $0.04$    & $0.58$ & $0.37$ & $0.01$ \\
            \end{tabular}\\[0.5em]
            {\scriptsize At $T{=}1$ cow and cat both look like zero; at $T{=}5$ the student can see that the
            teacher ranks cow well above cat. Adapted from Hinton, Vinyals \& Dean (2015).}
        \end{column}
    \end{columns}

\end{frame}

% --- SLIDE 6: DISTILLATION, THE LOSS ---
\begin{frame}{Distillation: The Loss, and What It Buys}

    \begin{itemize}
        \item The student minimises a mix of two cross-entropies:
              \[
                  L = \lambda\, \underbrace{\mathrm{CE}\big(y,\; q_S^{(T=1)}\big)}_{\text{match the labels}}
                  \;+\; (1-\lambda)\, T^2\, \underbrace{\mathrm{CE}\big(p_T^{(T)},\; q_S^{(T)}\big)}_{\text{match the teacher}}
              \]
        \item The soft term's gradients shrink as $1/T^2$, hence the $T^2$: without it, raising $T$ would
              quietly weaken the teacher's signal.
        \item \textbf{What it buys:} you choose the student's architecture, so it can be sized for the
              device. DistilBERT keeps 97\% of BERT's language-understanding performance with 40\% fewer
              parameters, and runs 60\% faster (Sanh et al., 2019).
        \item \textbf{What it costs:} a full training run, and the teacher's training data. The student also
              inherits the teacher's mistakes, including its weak slices.
    \end{itemize}

\end{frame}
